\chapter{Szakirodalmi háttér}

\section{Hálózatkutatás és közösségszerkezet}

A gráf mint modell evidens eszköz ismétlődő interakciók vizsgálatához. A videojátékos közösségek esetében különösen izgalmas kérdés, hogy a játékosok közötti ismétlődő interakciók mennyiben alkotnak felismerhető struktúrákat, szubközösségeket, hidakat és lokális mintázatokat. Girvan és Newman klasszikus munkája megmutatta, hogy a hálózat élei nem pusztán kapcsolódási elemek, hanem közösségszerkezetet leíró, eltérő fontosságú viszonyok is lehetnek \cite{girvannewman}. Fortunato átfogó áttekintése összefoglalja, miért központi probléma a közösségdetektálás a modern hálózatelemzésben \cite{fortunato2010}.

A jelen projektben a közösségfogalom nem öncélú. Azért fontos, mert a többször ismétlődő kapcsolatokra épülő játékoscsoportok értelmezhetőbbé teszik a nagyméretű hálózatot, segítenek elkülöníteni a valós kapcsolati magokat a zajtól és támaszt nyújtanak a későbbi útvonal- és analitikai számításoknak.

\section{Közösségdetektálás és modularitás}

A korai Python-oldali populációs klaszterezés a NetworkX \textit{greedy\_modularity\_communities} eljárását használta, amely a Clauset–Newman–Moore-féle mohó modularitásmaximalizálási megközelítésre épül \cite{clauset2004,networkx_greedy}. Ez tudatos kompromisszum volt: nem a leglátványosabb, hanem gyorsan magyarázható és reprodukálható megoldás. A jelenlegi futásidejű gráfállapotot a Rust motor önállóan építi fel a meccs-JSON-fájlokból és az SQLite-adatbázisból.

A modularitás-optimalizáció ismert korlátait sem hagyhatók figyelmen kívül: Fortunato és Barthelemy eredménye szerint a felbontási korlát miatt kisebb, de valós szerkezetek a globális optimum érdekében összeolvadhatnak \cite{fortunato_barthelemy2007}. A \textit{python\_population} klasztereket ezért nem szabad ontológiai ``igazi közösségként'' kezelni, inkább a választott gráf, súlymodell és heurisztika melletti legjobb, reprodukálható populációs partícióként.

\section{Asszortativitás}

Newman munkája az asszortatív keveredést olyan hálózati jelenségként írja le, amikor a kapcsolódó csomópontok hasonló vagy éppen eltérő tulajdonságúak \cite{newman2003}. A jelen projekt ezt a játékosok teljesítménymutatóira alkalmazza: nem az a kérdés, hogy a magas fokszámú csúcsok egymáshoz kapcsolódnak-e, hanem az, hogy a gráfban összekapcsolt játékosok hasonló \textit{opscore} vagy \textit{feedscore} értékeket hordoznak-e.

A dolgozat két gráfmódot különböztet meg: a \textit{social-path} csak a pozitív, szövetségesi támaszú kapcsolatokra épít, a \textit{battle-path} pedig az összes ismétlődő együtt- vagy egymás ellen játszást figyelembe veszi. Az asszortativitás mértéke e két módban eltér, és ez az eltérés számszerűen méri a kétféle gráfmód viselkedési különbségét.

\section{Játékos-teljesítmény és MOBA-analitika}

A League of Legends és az általánosabb MOBA-játékok esport-analitikája önálló kutatási területté vált. Mora-Cantallops és Sicilia átfogó irodalmi áttekintése megmutatja, hogy a MOBA-játékok egyedi kombinációt kínálnak verseny, együttműködés, játékosélmény és társas interakció között \cite{moracantallops2019}.

A játékos-teljesítmény modellezése Novak és munkatársai munkájában is megjelenik: profi League of Legends meccsekre kidolgozott teljesítménymodellek rámutatnak arra, hogy a domain-specifikus változók és a strukturált adatfeldolgozás fontosabbak, mint az általános statisztikai modellek \cite{novak2022}. Ez erősíti a jelen projekt szerepkör-érzékeny pontozási megközelítésének indokoltságát.

Sánchez-González és González-Bailón munkájában profi LoL-csapatok hálózati struktúráját kapcsolják össze a csapathatékonysággal \cite{sanchezgonzalezBailon2023}.

\section{Legrövidebb utak és heurisztikus keresés}

A klasszikus legrövidebb út problémát Dijkstra alapozta meg \cite{dijkstra1959}, míg az A* keresés formális alapját Hart, Nilsson és Raphael adták meg \cite{hart1968}. A Rust rétegben az egzakt A* implementáció ALT-szerű alsó becsléseket használ, amelyeket a Goldberg–Harrelson-féle megközelítés ihletett \cite{goldberg2005}. A cél a játékosháló nagyságrendjéhez illeszkedő, helyes és stabil futtatókörnyezet, nem új algoritmus, hanem meglévő, tesztelt megközelítések megbízható alkalmazása.

\section{Ágensorientált modellezés}

Az ágensorientált modellezés (ABM) olyan szimulációs paradigma, amelyben az egyéni szereplők lokális szabályok alapján cselekszenek, és a makroszintű minták e döntések összegzéseként jelennek meg. Bonabeau és munkatársai ezt a keretrendszert természetes és mesterséges rendszerekre egyaránt alkalmazták \cite{bonabeau2002}.

A dolgozat utolsó részét alkotó Tribal NeuroSim-kísérlet ebbe a paradigmába illeszkedik. Az ABM-irodalom hangsúlyozza: a szimulációból kapott eredményt a modell viselkedéseként kell értelmezni, nem a valóság közvetlen leképezéseként.
